\chapter{Computing time and GPU alternatives}
\label{ch:computation}

\section{Run times}
With the 8~cm mesh, one setting of 600~s took 15--90~hours on one CPU core of our lab computer, depending on
the window speeds (faster air needs smaller time steps). The runs only need the CPU; the \cotwo{} and the
training run on the GPU. \Cref{tab:speed} summarises tests on the slowest setting (N12, all windows at 2~m/s).

\begin{table}[h]
  \centering
  \caption{Speed tests on setting N12.}
  \label{tab:speed}
  \begin{tabular}{>{\raggedright}p{5cm}rp{6.3cm}}
    \toprule
    Change & Speed-up & Result\\
    \midrule
    4 cores instead of 1 & 2.7$\times$ & identical results; used when only few settings run\\
    450~s instead of 600~s & 1.33$\times$ & \cotwo{} 1.2--1.5\,\% different (N10: 3.1\,\%)\\
    300~s instead of 600~s & 2$\times$ & about 4\,\% different\\
    Courant number 2 instead of 1 & 1.8$\times$ & flow about 20\,\% different; rejected\\
    looser solver tolerances & 1.1$\times$ & not used\\
    \bottomrule
  \end{tabular}
\end{table}

For many settings it is most efficient to run one setting per core. The corrected dataset needed 28 new runs
of 150~s with a save every second; on the lab computer and additional machines they took less than two days.
Computing the \cotwo{} of one setting then takes 1.25--4~hours on the GPU, and the training of the 5 networks
about 1.5~hours.

\section{Lattice Boltzmann on the GPU}
To find out whether the simulations can be made much faster on the GPU, I wrote a lattice Boltzmann solver in
PyTorch (D3Q19, regularised BGK collision with a Smagorinsky model \citep{smagorinsky1963,kruger2017}). It uses
the same room, windows, door, columns and seating area, and computes the \cotwo{} in the same time loop. I first
looked at the existing PyTorch package Lettuce, but its outlets work only on a full side of the domain, not on a
door in part of a wall. I tested setting N10 on a 4~cm lattice (7 million nodes) for 600~s and compared it with
OpenFOAM:

\begin{itemize}
  \item it took 7.3~hours on the RTX A2000, which is not faster than OpenFOAM on 4 cores;
  \item the mean flow differs from OpenFOAM by 68\,\% and the \cotwo{} at 600~s by 71\,\% (paper's metric
        11.1\,\%, while my \cotwo{} solver on the OpenFOAM flow gives 4.1\,\%).
\end{itemize}

A closer look showed no error in the solver itself: the inflow speed at the window and the strength of the
fluctuations match OpenFOAM, and as much air leaves through the door as comes in. The difference comes from the turbulence model: in OpenFOAM the jet
from window W1 slows down to 0.27~m/s within 2~m, while in the lattice Boltzmann solution it still has 0.47~m/s
there. In addition, the mean over 150--600~s was not yet settled (17\,\% difference between the two halves). Since
neither result can be called correct without measurements, and the dataset should stay consistent, I did not use
this solver.

\section{Other software}
I checked the available software for running the same turbulence model (transient k-$\omega$ SST with a
transported scalar) on the GPU. \todo{add references for each}
\begin{itemize}
  \item \textbf{OpenFOAM with GPU linear solvers} (AmgX, PETSc4FOAM, Ginkgo): only the pressure solver runs on the
        GPU, so the total gain is limited.
  \item \textbf{OpenFOAM v2606}: the first official release with GPU offloading; still experimental, and it is
        not documented which solvers and turbulence models run on the GPU.
  \item \textbf{RapidCFD}: an OpenFOAM fork that runs fully on the GPU and contains \texttt{pimpleFoam},
        k-$\omega$ SST and our boundary conditions, but it is based on OpenFOAM 2.3.1 and CUDA 8 and could not be
        built with the current tools on our computer.
  \item \textbf{nekRS}: open source and GPU-based, with a related k-$\tau$ RANS model, but a different
        discretisation and a large effort to set up.
  \item \textbf{Lattice Boltzmann codes} (e.g.\ FluidX3D): fast, but LES instead of RANS (see above).
  \item \textbf{ANSYS Fluent}: its GPU solver supports k-$\omega$ SST, transient runs and species transport, so it
        could reproduce our setup. It needs a commercial or academic licence; the free student version seems to be
        limited to about 512,000 cells, less than our mesh. \todo{confirm the limit with official documentation}
\end{itemize}
\todo{result of a Fluent test if a licence is available; benchmark on DGX Spark and M5 Max}
